% ===========================================================================
%  Raport — Experiment: FP8 E4M3FN vs INT8 Weight-Only Quantization
%  ViT-Tiny pretrained pe ImageNet-1k validation
%  Tudose Alexandru · Lucrare de Licență 2026
% ===========================================================================
\documentclass[12pt,a4paper]{article}

\usepackage{fontspec}
\defaultfontfeatures{Ligatures=TeX}
\usepackage{polyglossia}
\setmainlanguage{romanian}
\usepackage[top=2.5cm, bottom=2.5cm, left=2.5cm, right=2.5cm]{geometry}
\usepackage{amsmath, amssymb}
\usepackage{graphicx}
\usepackage{float}
\usepackage{caption}
\usepackage{subcaption}
\captionsetup{font=small, labelfont=bf}
\usepackage{booktabs}
\usepackage{array}
\usepackage{xcolor}

\definecolor{fp32color}{HTML}{0072B2}
\definecolor{int8ptcolor}{HTML}{E69F00}
\definecolor{int8pccolor}{HTML}{D55E00}
\definecolor{fp8ptcolor}{HTML}{CC79A7}
\definecolor{fp8pccolor}{HTML}{009E73}

\usepackage[colorlinks=true, linkcolor=fp32color, citecolor=fp32color, urlcolor=fp32color]{hyperref}
\usepackage{microtype}
\usepackage{parskip}
\setlength{\parindent}{0pt}
\setlength{\parskip}{6pt}
\usepackage{enumitem}

\usepackage{fancyhdr}
\pagestyle{fancy}
\fancyhf{}
\lhead{\small Tudose Alexandru --- Experiment: FP8 E4M3FN vs INT8}
\rhead{\small 2026}
\cfoot{\thepage}
\renewcommand{\headrulewidth}{0.4pt}

% ===========================================================================
\begin{document}
% ===========================================================================

\begin{titlepage}
  \centering
  {\large Universitatea Transilvania din Brașov\\}
  {\small Facultatea de Matematică și Informatică\\}
  {\small Specializarea: Informatică Aplicată\\[2cm]}

  \rule{\linewidth}{0.4pt}\\[0.5cm]
  {\LARGE\bfseries Experiment: FP8 E4M3FN\\[0.3cm]}
  {\Large\bfseries vs INT8 Weight-Only Quantization\\[0.2cm]}
  {\large ViT-Tiny pe ImageNet-1k validation}\\[0.3cm]
  \rule{\linewidth}{0.4pt}\\[2cm]

  \begin{minipage}{0.45\textwidth}
    \textbf{Student:}\\
    Tudose Alexandru\\
    Grupa 10LF333\\
    An III, Informatică Aplicată
  \end{minipage}
  \hfill
  \begin{minipage}{0.45\textwidth}
    \raggedleft
    \textbf{Coordonator:}\\
    Conf.\ dr.\ ing.\ Honorius Gâlmeanu\\
    Informatică Aplicată
  \end{minipage}\\[3cm]
\end{titlepage}

\tableofcontents
\newpage

% ===========================================================================
\section{Introducere și Motivație}
% ===========================================================================

Studiul preliminar (Faza~0) a utilizat cuantizarea la \texttt{float8\_e4m3fn}
prin cast direct, fără stocare reală în FP8 --- ponderile erau reconvertite
imediat la FP32, deci nu exista reducere de memorie. Experimentul de față
implementează corect cuantizarea FP8, replicând structura din Faza~2
(INT8 cu stocare reală și dequantizare la forward pass), dar utilizând
formatul \texttt{float8\_e4m3fn} în loc de INT8.

Motivația vine dintr-o sugestie externă: dacă FP8 E4M3FN are o reprezentare
a numerelor în virgulă mobilă (exponent + mantisă), ar putea captura mai
bine distribuțiile non-uniforme ale ponderilor față de INT8 (reprezentare
uniformă pe întregi).

\subsection{FP8 E4M3FN --- format numeric}

Formatul \texttt{float8\_e4m3fn} (E4M3) utilizează:
\begin{itemize}
  \item 1 bit semn
  \item 4 biți exponent (bias = 7)
  \item 3 biți mantisă
  \item Domeniu: $\pm 448$, cu valori speciale (NaN, dar fără Inf)
  \item Densitate mai mare de valori reprezentabile în apropierea lui 0
\end{itemize}

Față de INT8 (256 valori uniform distribuite între $-128$ și $127$),
FP8 E4M3FN are o distribuție logaritmică --- mai mulți biți pentru
valorile mici, mai puțini pentru valorile mari. Ponderile rețelelor
neurale tind să aibă distribuții peakate în jurul lui 0, ceea ce
teoretic avantajează FP8.

\subsection{Diferența față de Faza 0}

\begin{table}[H]
  \centering
  \caption{Comparație abordări FP8}
  \begin{tabular}{lll}
    \toprule
    & \textbf{Faza 0 (preliminar)} & \textbf{Acest experiment} \\
    \midrule
    Stocare ponderi   & FP32 (cast temporar)   & \texttt{float8\_e4m3fn} (reală) \\
    Reducere memorie  & Nu                      & Da ($\approx 3.3\times$) \\
    Cuantizare        & Toate straturile        & Selectivă (skip norm, head) \\
    Dequantizare      & Imediată la cuantizare  & La fiecare forward pass \\
    Dataset           & CIFAR-10 (fine-tuned)   & ImageNet-1k (pretrained) \\
    \bottomrule
  \end{tabular}
  \label{tab:fp8_comparison}
\end{table}

% ===========================================================================
\section{Metodologie}
% ===========================================================================

\subsection{Implementare FP8Linear}

Clasa \texttt{FP8Linear} urmează același pattern ca \texttt{QuantizedLinear}
din Faza~2, cu adaptarea pentru formatul FP8:

\begin{equation}
  s = \frac{\text{FP8\_MAX}}{\max(|W|)}, \qquad
  q_{ij} = \text{clamp}(W_{ij} \cdot s,\; -448,\; 448).\text{to(float8)}
\end{equation}

La forward pass:
\begin{equation}
  \hat{W}_{ij} = q_{ij}.\text{to(float32)} \cdot \frac{1}{s}
\end{equation}

\textbf{Notă hardware:} Pe Apple M4, MPS nu suportă operații native cu
\texttt{float8\_e4m3fn}. Buffer-urile sunt stocate pe CPU, dequantizarea
se face pe CPU, iar tensorul rezultat este mutat pe MPS pentru matmul.
Pe A100/H100, dequantizarea și matmul-ul pot fi efectuate nativ în hardware.

\subsection{Configurații}

\begin{table}[H]
  \centering
  \caption{Cele 5 configurații evaluate}
  \begin{tabular}{llll}
    \toprule
    \textbf{Metodă} & \textbf{Format stocare} & \textbf{Scalare} & \textbf{Straturi} \\
    \midrule
    FP32       & float32  & ---           & 0 \\
    INT8-pt    & int8     & Per-tensor    & 48 \\
    INT8-pc    & int8     & Per-channel   & 48 \\
    FP8-pt     & float8\_e4m3fn & Per-tensor  & 48 \\
    FP8-pc     & float8\_e4m3fn & Per-channel & 48 \\
    \bottomrule
  \end{tabular}
  \label{tab:configs}
\end{table}

% ===========================================================================
\section{Rezultate}
% ===========================================================================

\subsection{Acuratețe și degradare}

\begin{table}[H]
  \centering
  \caption{Rezultate complete --- FP8 E4M3FN vs INT8}
  \begin{tabular}{lrrrr}
    \toprule
    \textbf{Metodă} & \textbf{Accuracy} & \textbf{$\Delta$ FP32} &
    \textbf{Loss} & \textbf{Mem (MB)} \\
    \midrule
    FP32       & 75.456\% & ---          & 0.9420 & 21.81 \\
    INT8-pt    & 75.134\% & $-0.322$~pp  & 0.9563 & \phantom{0}6.62 \\
    INT8-pc    & 75.424\% & $-0.032$~pp  & 0.9428 & \phantom{0}6.70 \\
    FP8-pt     & 75.258\% & $-0.198$~pp  & 0.9511 & \phantom{0}6.62 \\
    FP8-pc     & 75.356\% & $-0.100$~pp  & 0.9495 & \phantom{0}6.70 \\
    \bottomrule
  \end{tabular}
  \label{tab:results}
\end{table}

\begin{figure}[H]
  \centering
  \includegraphics[width=\textwidth]{../results/experiments/fp8_vs_int8/plots/01_degradation.png}
  \caption{Degradarea de acuratețe față de FP32 pentru toate configurațiile.
    FP8-pc ($-0.100$~pp) este mai bun decât INT8-pt ($-0.322$~pp) dar mai
    slab decât INT8-pc ($-0.032$~pp).}
  \label{fig:degradation}
\end{figure}

\begin{figure}[H]
  \centering
  \includegraphics[width=0.85\textwidth]{../results/experiments/fp8_vs_int8/plots/02_grouped.png}
  \caption{Comparație directă INT8 vs FP8 per granularitate. Per-tensor:
    FP8 este mai bun cu $+0.124$~pp față de INT8. Per-channel: INT8 este
    mai bun cu $+0.068$~pp față de FP8.}
  \label{fig:grouped}
\end{figure}

\begin{figure}[H]
  \centering
  \includegraphics[width=\textwidth]{../results/experiments/fp8_vs_int8/plots/00_summary.png}
  \caption{Sumar complet: acuratețe, memorie și latență pentru toate
    configurațiile. Memoria FP8 și INT8 este identică (aceeași dimensiune
    de stocare per parametru: 1 byte), latența FP8 este mai mare din cauza
    dequantizării pe CPU.}
  \label{fig:summary}
\end{figure}

\subsection{Eroarea de cuantizare (MSE)}

\begin{figure}[H]
  \centering
  \includegraphics[width=0.75\textwidth]{../results/experiments/fp8_vs_int8/plots/03_mse_comparison.png}
  \caption{MSE mediu per strat pentru INT8-pt, FP8-pt și INT8-pc. FP8-pt
    ($2.54 \times 10^{-6}$) are un MSE similar cu INT8-pt ($2.88 \times 10^{-6}$),
    ambele cu un ordin de mărime mai mari decât INT8-pc ($2.30 \times 10^{-7}$).}
  \label{fig:mse}
\end{figure}

\textbf{Observație cheie privind MSE-ul FP8:} Spre deosebire de INT8-pt
unde outlierii din block~7 aveau MSE de $3.5 \times 10^{-5}$ (de $\sim 12\times$
mai mare față de medie), în FP8-pt MSE-ul este uniform: maxim $4.11 \times 10^{-6}$
(\texttt{blocks.10.mlp.fc2}). FP8 rezolvă natural problema outlierilor prin
distribuția logaritmică a valorilor reprezentabile --- dar la un cost de
acuratețe globală mai mare față de INT8-pc.

\subsection{Tradeoff acuratețe--memorie}

\begin{figure}[H]
  \centering
  \includegraphics[width=0.85\textwidth]{../results/experiments/fp8_vs_int8/plots/04_tradeoff.png}
  \caption{Scatter plot acuratețe vs memorie. INT8-pt și FP8-pt au aceeași
    memorie ($6.62$~MB) dar acuratețe diferită. INT8-pc și FP8-pc au aceeași
    memorie ($6.70$~MB), INT8-pc fiind superior ca acuratețe.}
  \label{fig:tradeoff}
\end{figure}

% ===========================================================================
\section{Analiză și Interpretare}
% ===========================================================================

\subsection{Verdictul numeric}

\begin{table}[H]
  \centering
  \caption{Verdict: FP8 vs INT8 per granularitate}
  \begin{tabular}{lrrr}
    \toprule
    \textbf{Comparație} & \textbf{FP8} & \textbf{INT8} &
    \textbf{Diferență} \\
    \midrule
    Per-tensor & 75.258\% & 75.134\% & FP8 mai bun cu $+0.124$~pp \\
    Per-channel & 75.356\% & 75.424\% & INT8 mai bun cu $+0.068$~pp \\
    \bottomrule
  \end{tabular}
  \label{tab:verdict}
\end{table}

Rezultatele sunt nuanțate:
\begin{itemize}
  \item \textbf{Per-tensor:} FP8 câștigă ($+0.124$~pp). Distribuția
    logaritmică a FP8 compensează parțial problema outlierilor din block~7,
    pe care INT8 per-tensor o gestionează prost.
  \item \textbf{Per-channel:} INT8 câștigă ($+0.068$~pp). Per-channel
    rezolvă deja problema outlierilor prin scale-uri individuale pe canal,
    iar INT8 oferă mai multă precizie uniformă decât FP8 în această
    configurație.
\end{itemize}

\subsection{De ce FP8-pt $>$ INT8-pt?}

INT8-pt folosește o singură scalare per matrice. Ponderile din block~7
(\texttt{mlp.fc1}, \texttt{mlp.fc2}) au outlieri cu valori
$\sim 10\times$ mai mari decât restul, ceea ce face ca scala globală să
fie prea mare pentru valorile mici --- erori de cuantizare ridicate
pentru majoritatea parametrilor.

FP8 E4M3FN alocă automat mai mulți biți valorilor mici (distribuție
logaritmică). Deși scala e tot globală (per-tensor), formatul FP8
capturează mai bine distribuțiile cu coadă lungă, explicând MSE-ul
mai uniform și acuratețea mai bună față de INT8-pt.

\subsection{De ce INT8-pc $>$ FP8-pc?}

Per-channel elimină problema outlierilor prin scale individuale ---
INT8-pc a redus MSE-ul cu $12.52\times$ față de INT8-pt (Faza~3).
Odată rezolvată problema outlierilor, INT8 devine superior: 256 de
valori uniform distribuite per canal oferă mai multă precizie pe
distribuțiile gaussiene ale ponderilor decât cei 256 de pași neuniformi
ai FP8.

\subsection{Latența pe Apple M4}

\begin{table}[H]
  \centering
  \caption{Latență per configurație (ms/batch)}
  \begin{tabular}{lr}
    \toprule
    \textbf{Metodă} & \textbf{Latență (ms/batch)} \\
    \midrule
    FP32    & 166.5 \\
    INT8-pt & 210.5 \\
    INT8-pc & 202.8 \\
    FP8-pt  & 250.7 \\
    FP8-pc  & 274.8 \\
    \bottomrule
  \end{tabular}
  \label{tab:latency}
\end{table}

FP8 este semnificativ mai lent pe Apple M4: FP8-pc ($274.8$~ms) față de
INT8-pc ($202.8$~ms), un overhead de $35\%$. Motivul: buffer-urile FP8
sunt stocate pe CPU (MPS nu suportă \texttt{float8\_e4m3fn}), dequantizarea
se face pe CPU, iar transferul CPU→MPS adaugă latență la fiecare strat
și la fiecare batch. Pe hardware cu suport FP8 nativ (H100 Hopper),
situația s-ar inversa --- FP8 ar putea fi mai rapid decât INT8.

% ===========================================================================
\section{Concluzii}
% ===========================================================================

\begin{enumerate}
  \item \textbf{FP8-pt bate INT8-pt cu $+0.124$~pp:} Distribuția
    logaritmică a FP8 gestionează mai bine outlierii față de scalarea
    uniformă INT8 per-tensor, fără a necesita scale per-canal.

  \item \textbf{INT8-pc rămâne cel mai bun format:} $75.424\%$,
    $-0.032$~pp față de FP32, cu $3.26\times$ reducere de memorie.
    FP8-pc ($75.356\%$, $-0.100$~pp) este inferior.

  \item \textbf{FP8 uniformizează MSE-ul:} Spre deosebire de INT8-pt
    unde block~7 avea MSE de $3.5 \times 10^{-5}$, în FP8-pt MSE-ul
    maxim este $4.1 \times 10^{-6}$ --- o distribuție mult mai uniformă.

  \item \textbf{Memoria identică:} FP8 și INT8 ocupă același spațiu
    ($6.62$--$6.70$~MB), deoarece ambele stochează 1 byte per parametru.

  \item \textbf{FP8 necesită hardware dedicat:} Pe Apple M4, latența
    FP8 este $35\%$ mai mare față de INT8. Beneficiile reale ale FP8
    (speedup hardware + precizie îmbunătățită față de INT8-pt) se
    manifestă pe GPU-uri Hopper (H100) cu suport nativ pentru FP8 GEMM.

  \item \textbf{Recomandare generală:} Fără hardware FP8 nativ,
    INT8-pc rămâne soluția optimă. Cu hardware H100, FP8-pt devine
    competitiv --- mai simplu de implementat decât per-channel și
    cu acuratețe mai bună decât INT8-pt.
\end{enumerate}

% ===========================================================================
\begin{thebibliography}{9}
% ===========================================================================

\bibitem{micikevicius2022fp8}
P. Micikevicius, D. Stosic, N. Burgess et al.,
``FP8 Formats for Deep Learning,''
2022. \href{https://arxiv.org/abs/2209.05433}{arXiv:2209.05433}.

\bibitem{dosovitskiy2021vit}
A. Dosovitskiy et al.,
``An Image is Worth 16x16 Words: Transformers for Image Recognition at Scale,''
\emph{ICLR}, 2021. \href{https://arxiv.org/abs/2010.11929}{arXiv:2010.11929}.

\bibitem{nagel2021whitepaper}
M. Nagel et al.,
``A White Paper on Neural Network Quantization,''
2021. \href{https://arxiv.org/abs/2106.08295}{arXiv:2106.08295}.

\bibitem{dettmers2022llmint8}
T. Dettmers et al.,
``LLM.int8(): 8-bit Matrix Multiplication for Transformers at Scale,''
\emph{NeurIPS}, 2022. \href{https://arxiv.org/abs/2208.07339}{arXiv:2208.07339}.

\bibitem{nvidia_h100}
NVIDIA Corporation,
``H100 Tensor Core GPU Architecture,''
2022. \url{https://resources.nvidia.com/en-us-tensor-core}.

\end{thebibliography}

\end{document}
